\documentclass[11pt,a4paper]{article}

\usepackage[utf8]{inputenc}
\usepackage[T1]{fontenc}
\usepackage{lmodern}
\usepackage[brazil]{babel}
\usepackage[a4paper,margin=2.5cm]{geometry}
\usepackage{xcolor}
\definecolor{captagov}{HTML}{B8791E}
\definecolor{captagovdark}{HTML}{7A5013}
\usepackage[colorlinks=true,linkcolor=captagovdark,urlcolor=captagovdark,citecolor=captagovdark]{hyperref}
\usepackage{titlesec}
\usepackage{booktabs}
\usepackage{longtable}
\usepackage{array}
\usepackage{enumitem}
\usepackage{parskip}
\usepackage{fancyhdr}

\titleformat{\section}{\normalfont\Large\bfseries\color{captagovdark}}{\thesection}{0.6em}{}
\titleformat{\subsection}{\normalfont\large\bfseries\color{captagovdark}}{\thesubsection}{0.6em}{}
\titleformat{\subsubsection}{\normalfont\normalsize\bfseries}{\thesubsubsection}{0.6em}{}

\pagestyle{fancy}
\fancyhf{}
\renewcommand{\headrulewidth}{0.4pt}
\fancyhead[L]{\small CaptaGov --- Relatório Técnico e Roadmap}
\fancyhead[R]{\small \thepage}

\newcommand{\codigo}[1]{\texttt{\small #1}}

\title{\textbf{\color{captagov}CaptaGov}\\[0.3em]\Large Relatório Técnico e Roadmap}
\author{Matheus Assis de Oliveira}
\date{22 de agosto de 2026}

\begin{document}
\maketitle
\thispagestyle{fancy}

\begin{abstract}
\noindent Este documento descreve o estado técnico atual do protótipo \textbf{CaptaGov} --- uma
ferramenta que cruza programa de repasse federal aberto com o histórico de captação de cada
município, apontando oportunidade não usada --- e propõe quatro extensões técnicas para
evoluir a ferramenta dentro de uma família maior de produtos voltados à gestão pública
estadual e municipal (que já inclui \textbf{CityPro}, para licenciamento urbano, e
\textbf{FLORA}, para governança ambiental).
\end{abstract}

\section{Contexto}

O CaptaGov nasceu de uma pergunta em aberto: quais ferramentas fariam sentido para atender
estados e prefeituras. Duas referências de mercado já existentes nesse espaço foram
usadas como ponto de partida:

\begin{itemize}
  \item \textbf{CityPro} --- plataforma de licenciamento urbano com leitura automatizada de
    projeto via BIM, focada em reduzir o tempo de aprovação de obra.
  \item \textbf{FLORA} --- plataforma de gestão ambiental pública, integrando licenciamento
    ambiental, fiscalização, monitoramento territorial e atendimento a denúncia.
\end{itemize}

Ambas resolvem o lado da \emph{regulação} (aprovar, fiscalizar, licenciar). Nenhuma delas
ataca o lado da \emph{captação de recurso} --- e é exatamente aí que a maioria dos municípios
pequenos brasileiros mais perde dinheiro: falta de corpo técnico para monitorar edital e
programa de repasse federal aberto, e falta de visibilidade sobre o que já foi captado. O
CaptaGov nasce para preencher essa lacuna, como terceiro pilar de uma mesma família de
produtos.

\section{CaptaGov --- Visão Geral}

\subsection{O problema}

Prefeituras de pequeno e médio porte raramente têm equipe dedicada a monitorar oportunidade
de captação de recurso federal. O resultado é previsível: programa de repasse fica aberto por
meses, elegível para aquele município, e ninguém submete proposta a tempo --- dinheiro que
ficou disponível e não foi buscado.

\subsection{A proposta de valor}

O CaptaGov automatiza três passos que hoje são feitos manualmente (quando são feitos):

\begin{enumerate}
  \item \textbf{Descobrir} quais programas de repasse federal estão abertos agora para
    candidatura direta (sem depender de indicação de parlamentar).
  \item \textbf{Cruzar} essa lista com o histórico de captação de cada município, calculando
    o \emph{gap}: programa que o município é elegível e ainda não usou.
  \item \textbf{Gerar} um rascunho de plano de trabalho de partida, já preenchido com os
    campos oficiais do programa (objetivo, problema, público-alvo, resultado esperado,
    condicionantes, documentos exigidos).
\end{enumerate}

\section{Arquitetura Técnica}

\subsection{Fontes de dado}

Todo o dado do CaptaGov vem de \textbf{APIs públicas do governo federal, sem autenticação e
sem custo}. Duas fontes distintas são combinadas:

\begin{description}[leftmargin=1.2cm,style=nextline]
  \item[API de dados abertos] \codigo{api-publica.transferegov.gestao.gov.br/parcerias} ---
    módulo ``Parcerias'' do Transferegov.br (sucessor do SICONV). Expõe os endpoints
    \codigo{/programa} (programa de repasse cadastrado, com situação, tipo de beneficiário,
    objetivo, problema, público-alvo) e \codigo{/proposta} (proposta já enviada por um
    município, com valor e situação). Filtra por \codigo{cd\_ibge\_recebedor} (código IBGE do
    município).
  \item[Portal do Transferegov (detalhe rico)] \codigo{parcerias.transferegov.sistema.gov.br}
    --- o mesmo sistema que os servidores públicos usam para gerir os programas. Por trás da
    interface web existe uma API JSON pública (sem login),
    \codigo{/ep/api/atos-prep/programa/\{id\}}, que expõe campos não disponíveis na API de
    dados abertos: requisitos e condicionantes para recepção de recurso, documentos e anexos
    do programa (ex.\ portaria, resolução), e a janela real de captação
    (\codigo{dtaInicioRecebPropEspfic} / \codigo{dtaFimRecebPropEspfic}). O mesmo \codigo{id}
    de programa vale nas duas APIs, o que permite cruzar os dois sem trabalho extra.
  \item[IBGE] Lista de municípios e código IBGE via
    \codigo{servicodados.ibge.gov.br/api/v1/localidades}; coordenadas de cada município
    calculadas a partir da malha municipal
    (\codigo{.../api/v3/malhas/estados/\{uf\}?intrarregiao=municipio}), tomando a média dos
    vértices do polígono como centroide aproximado --- suficiente para posicionar um marcador
    no mapa, sem precisar de geocoding cidade a cidade.
\end{description}

\subsection{Pipeline de dados}

O módulo \codigo{data\_fetch.py} concentra toda a lógica de acesso a dado:

\begin{itemize}
  \item \textbf{Cache local em disco}, com \emph{time-to-live} de 24 horas por arquivo JSON
    --- dado de governo não muda a cada minuto, e isso evita bater a API repetidamente durante
    desenvolvimento ou geração do painel.
  \item \codigo{get\_programas\_abertos()} --- filtra os programas com situação
    \emph{Disponibilizado} e qualificação de beneficiário \emph{Espontâneo} ou
    \emph{Específico} (ou seja, candidatura direta, sem depender de indicação de emenda
    parlamentar).
  \item \codigo{get\_propostas\_municipio(cd\_ibge)} --- histórico de propostas já enviadas
    por um município.
  \item \codigo{calcular\_gap(...)} --- programa aberto que o município ainda não tem proposta
    vinculada.
  \item \codigo{get\_programa\_detalhe(id\_programa)} --- busca o detalhe rico no portal
    (requisitos, anexos, janela de captação real), com o mesmo padrão de cache.
  \item \codigo{gerar\_minuta(...)} --- monta o rascunho de plano de trabalho combinando os
    campos das duas fontes.
\end{itemize}

\subsection{Duas superfícies de entrega}

O mesmo pipeline de dado alimenta duas formas de uso, propositalmente distintas:

\begin{enumerate}
  \item \textbf{Painel estático autocontido} (\codigo{gerar\_painel.py}
    $\rightarrow$ \codigo{docs/index.html}) --- roda uma vez, embute todo o dado como JSON
    dentro de um único arquivo HTML. Não precisa de servidor, banco de dado ou instalação para
    ser visualizado: abre por duplo clique em qualquer computador, e pode ser publicado direto
    num GitHub Pages. É a superfície pensada para \emph{mostrar e compartilhar} o protótipo.
    Tem como custo o dado ficar ``congelado'' no momento em que o gerador rodou --- para
    atualizar, roda de novo.
  \item \textbf{Servidor Flask} (\codigo{app.py}) --- serve o mesmo dado sempre ao vivo,
    buscando a API a cada visita nova (respeitando o cache de 24h). Pensado para uso local /
    desenvolvimento, onde dado sempre atualizado importa mais que portabilidade.
\end{enumerate}

\subsection{Interface do painel}

O painel estático foi desenhado no mesmo padrão visual de outro produto do autor
(\emph{atf-georadar}, ferramenta de inteligência imobiliária): tudo em uma única página,
sem \emph{build step}, usando bibliotecas via CDN (\emph{Chart.js} para os gráficos,
\emph{Leaflet} para o mapa) e dado embutido diretamente no HTML gerado. Recursos:

\begin{itemize}
  \item \textbf{KPIs} --- total captado historicamente, número de programas abertos, total de
    oportunidade não usada, número de cidades cobertas.
  \item \textbf{Gráficos} --- valor histórico captado por cidade e oportunidade aberta não
    usada por cidade, ambos em barra, ordenados.
  \item \textbf{Mapa} --- um marcador por município (tamanho proporcional à oportunidade
    aberta não usada), clicável, filtra a tabela ao clicar.
  \item \textbf{Filtro cruzado} --- cidade, órgão repassador e busca por texto no nome do
    programa, atualizando tabela, gráfico e histórico juntos.
  \item \textbf{Ficha / minuta imprimível} --- modal com o rascunho de plano de trabalho, link
    direto para o programa no portal oficial, e botão de impressão / salvar em PDF.
  \item \textbf{Exportar CSV} --- baixa a seleção filtrada atual.
  \item \textbf{Sistema de temas} --- três paletas trocáveis (dourado, verde, azul),
    selecionadas por um controle no canto do cabeçalho e persistidas em
    \codigo{localStorage}. O dourado é o tema padrão do CaptaGov; verde e azul dialogam
    respectivamente com a identidade visual do FLORA e com um tom institucional mais neutro,
    mantendo os três produtos da família reconhecíveis lado a lado (fundo bem escuro, cantos
    bem arredondados, badge em formato de pílula, marca de canto), sem repetir a cor de nenhum
    dos outros dois produtos.
\end{itemize}

\subsection{Estado atual dos dados}

Na geração mais recente do painel, cobrindo os 79 municípios de Mato Grosso do Sul:

\begin{center}
\begin{tabular}{@{}lr@{}}
\toprule
\textbf{Indicador} & \textbf{Valor} \\
\midrule
Municípios cobertos & 79 (todos os do MS) \\
Programas federais abertos para candidatura direta & 18 \\
Oportunidades (cidade × programa não usado) & 1.385 \\
Valor histórico de captação mapeado & R\$ 839.437.599 \\
Propostas históricas indexadas & 1.341 \\
\bottomrule
\end{tabular}
\end{center}

\section{Limitações Conhecidas}

Transparência técnica sobre o que o protótipo \emph{não} faz ainda:

\begin{itemize}
  \item A minuta gerada é um rascunho de texto simples a partir dos campos do programa --- não
    é o formulário oficial de nenhum órgão repassador, que varia por ministério.
  \item Só o módulo ``Parcerias'' do Transferegov foi explorado. Existem outros três módulos
    (Especiais, Fundo a Fundo, Discricionárias e Legais) ainda não integrados.
  \item O painel estático publicado não se atualiza sozinho; reflete o momento em que o
    gerador rodou pela última vez.
  \item Mapa e gráfico dependem de carregar Chart.js e Leaflet via CDN --- funcionam offline
    apenas depois do primeiro carregamento com internet.
\end{itemize}

\section{Roadmap Técnico --- Próximas Integrações}

Quatro extensões foram identificadas como próximo passo natural, cada uma com uma trava de
acesso a dado diferente --- distinção que muda completamente o que é possível prototipar
agora versus o que só vira produto real depois de uma relação comercial com o cliente.

\subsection{Prestação de contas automática}

\textbf{O que resolve:} depois que o recurso é captado, a prefeitura precisa prestar contas
do uso --- processo que hoje assusta gestor pequeno pelo risco de glosa e inabilitação junto
ao Tribunal de Contas. É a continuação natural do CaptaGov: hoje ele só ajuda a captar, essa
extensão ajuda a não perder o que já foi captado.

\textbf{Fonte de dado:} mesma API pública já em uso, endpoint
\codigo{/ep/api/atos-prep/parceria/editar/\{id\}} (já mapeado durante o desenvolvimento do
CaptaGov), que expõe plano de trabalho, metas padrão e dado orçamentário da parceria já
formalizada.

\textbf{Como implementaria:} novo módulo seguindo o mesmo padrão de
\codigo{data\_fetch.py} (cache de 24h, mesma convenção de nomes), cruzando meta planejada no
plano de trabalho contra execução relatada, gerando um documento de prestação de contas
estruturado a partir do gap entre os dois.

\textbf{Trava de acesso:} nenhuma --- mesmo dado público, mesmo pipeline. É a única das
quatro extensões que pode evoluir hoje sem depender de ninguém liberar acesso a nada.

\subsection{Cadastro imobiliário automatizado (IPTU)}

\textbf{O que resolve:} baixa arrecadação própria por cadastro imobiliário desatualizado ---
construção nova ou ampliação não declarada ao município. Ferramenta que aumenta receita sem
alterar alíquota, o que politicamente é muito mais fácil de vender a um prefeito do que
aumento de imposto.

\textbf{Fonte de dado:} reaproveita o motor geoespacial já construído no \emph{atf-georadar}
(ponto-em-polígono por \emph{ray-casting} próprio, malha de lote oficial via CTM,
zoneamento via servidor ArcGIS da prefeitura), somado a imagem de satélite (Sentinel-2,
gratuita, ou aérea da prefeitura quando disponível) para detectar área construída.

\textbf{Como implementaria:} pipeline de detecção de contorno construído sobre a imagem de
satélite, comparado ao lote da malha CTM; cruzamento contra o cadastro de IPTU para apontar
divergência entre área construída real e área declarada.

\textbf{Trava de acesso:} o cadastro de IPTU é um sistema interno da prefeitura (ex.\ SIAT,
e-Cidade), sem API pública. Um protótipo funciona só com dado público --- mostra o método
(``este imóvel tem contorno construído que não bate com a média da região''), mas sem cruzar
contra o cadastro real. O produto de verdade só existe depois que a prefeitura vira cliente e
libera acesso ao próprio cadastro.

\subsection{Fiscalização de obra irregular}

\textbf{O que resolve:} fecha o ciclo com o CityPro --- de nada adianta um licenciamento
urbano eficiente se obra clandestina (sem alvará ou fora do projeto aprovado) segue sem
fiscalização.

\textbf{Fonte de dado:} alvará emitido (dado que só existe dentro do banco do CityPro,
produto fechado de terceiro) cruzado com imagem de satélite temporal, comparando duas datas
para detectar mudança física no lote.

\textbf{Como implementaria:} pipeline de detecção de mudança (\emph{change detection}) entre
duas capturas de imagem do mesmo lote; sinaliza lote com construção nova sem alvará
correspondente, ou com área construída maior que a aprovada.

\textbf{Trava de acesso:} igual ao item anterior --- só integra de fato com acesso ao banco
de dado do CityPro, o que depende de parceria comercial com quem desenvolveu o produto, não
de scraping de dado público. Protótipo demonstrativo usaria alvará simulado.

\subsection{Painel de indicadores municipais unificado}

\textbf{O que resolve:} hoje cada ferramenta da família fala de um recorte isolado --- CityPro
fala de licenciamento, FLORA de meio ambiente, CaptaGov de captação. Um painel único, cruzando
indicador de saúde, educação, meio ambiente e finanças públicas do mesmo município, é o
próximo passo depois que os produtos pontuais provarem valor individualmente.

\textbf{Fonte de dado:} múltiplas APIs públicas --- IBGE SIDRA (indicador socioeconômico),
DataSUS / e-Gestor AB (saúde), INEP / QEdu (educação), além do próprio Transferegov já
integrado.

\textbf{Como implementaria:} um conector por fonte, cada um normalizando pro mesmo identificador de
município (código IBGE), consolidando num único painel comparável.

\textbf{Trava de acesso:} nenhuma trava de acesso, mas custo de manutenção mais alto --- cada
fonte tem formato, cadência de atualização e nomenclatura próprios; é o item de maior esforço
de engenharia entre os quatro, e o de proposta de valor mais genérica isoladamente.

\subsection{Resumo comparativo}

\begin{center}
\begin{tabular}{@{}p{4.2cm}p{3.6cm}p{3.6cm}p{2.6cm}@{}}
\toprule
\textbf{Extensão} & \textbf{Fonte de dado} & \textbf{Trava de acesso} & \textbf{Esforço} \\
\midrule
Prestação de contas & API pública Transferegov & Nenhuma & Baixo \\
Cadastro imobiliário (IPTU) & Satélite + malha CTM & Cadastro interno da prefeitura & Médio \\
Fiscalização de obra irregular & Satélite temporal + CityPro & Banco de dado do CityPro & Médio \\
Painel unificado & IBGE, DataSUS, INEP, Transferegov & Nenhuma & Alto \\
\bottomrule
\end{tabular}
\end{center}

\section{Conclusão e Próximos Passos}

O CaptaGov já demonstra, com dado real e sem custo de API, o método central da família de
produtos: cruzar dado público disperso para revelar oportunidade que o gestor público não
enxerga sozinho. A extensão de \textbf{prestação de contas} é o passo mais natural a seguir
--- zero trava de acesso, mesmo pipeline, resolve a dor que mais preocupa prefeito pequeno
depois de captar recurso.

As duas extensões que dependem de dado interno da prefeitura ou do CityPro (\textbf{cadastro
de IPTU} e \textbf{fiscalização de obra irregular}) são, ao mesmo tempo, as de maior potencial
de receita e as que exigem uma conversa comercial antes de virarem produto de verdade --- o
protótipo com dado público serve para abrir essa conversa, não para substituí-la.

\vspace{1em}
\noindent\textbf{Link do painel publicado:}
\url{https://matheusassiso.github.io/captagov-prototype/}

\end{document}
